\documentclass[11pt]{article}
\usepackage[letterpaper,margin=0.88in]{geometry}
\usepackage[T1]{fontenc}
\usepackage{lmodern,microtype,amsmath,amssymb,amsthm,mathtools}
\usepackage{booktabs,longtable,array,enumitem,xcolor,xurl,hyperref,fancyhdr}
\hypersetup{colorlinks=true,linkcolor=blue!45!black,urlcolor=blue!45!black,citecolor=blue!45!black,
 pdftitle={Weaker congruences and even Pell indices},pdfauthor={OpenAI ChatGPT; research note for Vladimir Reshetnikov}}
\setlength{\parskip}{4pt}\setlength{\parindent}{0pt}
\setlist{itemsep=3pt,topsep=4pt}
\pagestyle{fancy}\fancyhf{}\fancyhead[L]{\small Explicit Diophantine consistency equations}
\fancyhead[R]{\small 24 October 2026}\fancyfoot[C]{\thepage}
\setlength{\headheight}{14pt}
\newcommand{\Z}{\mathbb Z}\newcommand{\N}{\mathbb N}
\newcommand{\Con}{\operatorname{Con}}\newcommand{\ZFC}{\mathsf{ZFC}}
\newcommand{\NN}{\mathcal N_4}\newcommand{\PP}{\mathcal P}
\newcommand{\ord}{\operatorname{ord}}\newcommand{\sgn}{\operatorname{sgn}}
\newcommand{\Pair}{\operatorname{Pair}}\newcommand{\bits}{\operatorname{bits}}
\newcommand{\file}[1]{\nolinkurl{#1}}
\newtheorem{theorem}{Theorem}[section]
\newtheorem{lemma}[theorem]{Lemma}\newtheorem{proposition}[theorem]{Proposition}
\newtheorem{corollary}[theorem]{Corollary}
\theoremstyle{definition}\newtheorem{definition}[theorem]{Definition}
\theoremstyle{remark}\newtheorem{remark}[theorem]{Remark}
\allowdisplaybreaks
\begin{document}
\begin{center}
{\LARGE\bfseries Weaker congruences and even Pell indices}\\[7pt]
{\Large Smaller explicit equations for $\Con(\ZFC)$}\\[10pt]
{\normalsize A research package for Vladimir Reshetnikov's ProveIt project}\\[3pt]
{\normalsize 24 October 2026}
\end{center}
\begin{abstract}
We give two concrete improvements to the explicit Diophantine consistency
equations in a pinned ProveIt snapshot. First, we weaken a divisibility
carrier, use the actual coefficient size to resolve its fifth-power
congruence, and combine two congruences in one homogeneous intersective
polynomial. The resulting equations have twelve unrestricted integer
unknowns and exact degrees $768,784,800,816,1040$, with respectively
$161,154,147,144,140$ binary arithmetic operations. The prior comparable
points had degrees $1200$ and $1456$ with $145$ and $141$ operations.
Second, in the direct C222 construction, we combine signs for repeated
bounded intervals and replace power requests $m=n^k$ by $m^2=n^{2k}$.
Even Pell indices eliminate twelve index-sign predicates. The complete
result is a quartic in $566$ unrestricted integers, with $6431$ monomials
and coefficient height $2445790974704$, compared with $604$ variables and
$9676$ monomials at the same height in the pinned parent. Every equation,
constant payload, projection certificate and verification program is
supplied. The new reductions have ordinary unbounded proofs; their
consistency interpretation composes with the explicitly cited existing
source/calculus calibration. No new Lean or Rocq end-to-end proof,
unconditional consistency assertion, global optimum, or literature-wide
priority claim is made.
\end{abstract}

\section{What is constructed, and what is improved}

For each final polynomial $P$ in this package, the reductions below and
the cited source calibration give
\begin{equation}\label{eq:target}
 \ZFC\vdash
 \left(\neg\exists\mathbf x\in\Z^n\;P(\mathbf x)=0\right)
 \ \longleftrightarrow\ \Con(\ZFC).
\end{equation}
Here consistency is the usual arithmetic assertion that there is no finite
syntactic proof of a contradiction in the fixed calculus. We prove
equivalence of \emph{existence} of polynomial zeros to finite refutability.
We do not prove that the displayed polynomials have no zeros.

The comparison is to ProveIt commit
\file{84e16e43092614bae262ca59651cb9d77e31a255}, including its active research
and incoming reports, rather than only the older collection README
\cite{repo,latest}. ``Simplest'' is not a total order: variable count,
degree, operation count, coefficient size and proof complexity can conflict.
The following are actual, fully materialized points.

\begin{center}
\begin{tabular}{lrrrr}
\toprule
Construction & Integer variables & Degree & Operations & Status\\
\midrule
Prior scalar carrier, rational &12&1200&145& inherited\\
Prior scalar carrier, integral &12&1456&141& inherited\\
Weak carrier, integral &12&1040&140& new\\
Weak carrier, rational &12&816&144& new\\
Centered weak carrier &12&800&147& new\\
Mixed-modulus carrier, integral &12&784&154& new\\
Mixed-modulus carrier, rational &12&768&161& new\\
\bottomrule
\end{tabular}
\end{center}
Operations count every live binary addition, subtraction and multiplication,
including squares and scalar products. Named polynomial values add no
unknowns. All six large coefficient values are unchanged in these comparisons.
The $140/1040$ and $144/816$ points improve both degree and operations
relative to the corresponding old points. The lower-degree points make
explicit cost/degree tradeoffs.

\begin{center}
\begin{tabular}{lrrrr}
\toprule
Direct quartic & Variables & Quadratic rows & Monomials & Height\\
\midrule
Pinned parent &604&372&9676&$2445790974704$\\
Seven sign combinations &589&371&8157&$2445790974704$\\
Even-index compiler &566&372&6431&$2445790974704$\\
\bottomrule
\end{tabular}
\end{center}
The last reduction removes $38$ variables and $3245$ monomials from the
matched direct parent, approximately $6.29\%$ and $33.54\%$, respectively.
The direct branch uses C222, not K215 or Jones's universal polynomial.
The twelve-variable branch retains the repository's existing native source.

Ten-integer high-degree constructions already exist in the same repository;
Sun's general ten-integer theorem is also a relevant benchmark
\cite{sun}. The present twelve-variable equations do not improve the
least known variable count. Their concrete gain is the degree/cost frontier
within a completely instantiated family.

\section{The complete twelve-variable equation}
\label{sec:equation}

The only quantified unknowns are
\[
 (j,C,z,h,r,v,w,T,c,n,t,s)\in\Z^{12}.
\]
All other symbols below are polynomial abbreviations. The six fixed
constants are $\beta,K_0,M_C,M_F,A_c,B_c$ and are specified in Section~\ref{sec:constants}.
Define the source front by
\begin{align}
 J&=6j+1,& q&=\beta J+1,& W&=12w+4,\\
 X&=qW,&Y&=qv,& a&=(X+1)Y,\\
 H&=4a+3,& \Delta&=a^2+H=(a+2)^2-1,\label{eq:base}\\
 F&=(K_0+W)C-z(q-1),\nonumber\\
 R&=\bigl(q(q-F)-C\bigr)(q^2-1)+(M_C+qM_F)J,\label{eq:R}\\
 k&=R+1+hXY,& \eta&=c-Yk,&\zeta&=k-\eta.
\end{align}
The strict source guards are encoded through
\begin{align}
 G&=(A_cF)(B_cC)-(q-A_cF-B_cC)^2,\\
 E&=\eta\zeta,\\
 B_0&=GE-(t^2-G-E+1)^2,\\
 D_g&=(3B_0-4)s^2+1.\label{eq:guard}
\end{align}
The Pell registers are
\begin{align}
 I&=src^2,& D_f&=1+\Delta I^2,& Q&=\Delta^2I^2,\\
 \delta&=Q-1,&V&=D_f(c-R+4IT)-c,& C_y&=QV^2-1,\\
 \ell&=XY^2,&D_\tau&=1+\ell(\ell+1)k^2.
\end{align}
Finally define the new carrier:
\begin{align}
 K_*&=X(2X-3c)-2,& M&=2q^2,&D&=MH,\\
 C_*&=MK_*+H(v-q^2),& z_*&=DT+C_*,\\
 A_*&=(z_*^2+3D^2)(z_*^3-19D^3),& S&=srM^5H.\label{eq:carrier}
\end{align}

\begin{definition}[Four-root norm]
Let
\begin{equation}\label{eq:norm}
 \NN(x;a,b,e,d)=
 \prod_{\varepsilon\in\{\pm1\}^4}
 \bigl(x-\varepsilon_1\sqrt a-\varepsilon_2\sqrt b
          -\varepsilon_3\sqrt e-\varepsilon_4\sqrt d\bigr).
\end{equation}
This is an ordinary integer polynomial in $x,a,b,e,d$: changing the sign
of any formal square root permutes the factors. Its degree in $x$ is $16$,
and its degree in each radicand is $8$.
\end{definition}

Our main equation is
\begin{equation}\label{eq:768}
\boxed{\begin{aligned}
 P_{768}={}&\delta^8\NN\bigl(S(n+V)-A_*;\\
 &9S^2D_f,\;81S^2D_\tau,\;729S^2D_g,\;S^2C_y/\delta\bigr)=0.
\end{aligned}}
\end{equation}
The division in this display is only a compact presentation. Expanding
the polynomial norm and multiplying by $\delta^8$ clears every denominator.
The supplied $161$-operation DAG contains no division, radical, exponent
with a variable exponent, or existential subroutine. Appendix~\ref{app:norm}
gives a division-free evaluator.

The companion equation is
\begin{equation}\label{eq:784}
\boxed{\begin{aligned}
 P_{784}={}&\NN\bigl(S(n+\delta V)-A_*;\\
 &9S^2D_f,\;81S^2D_\tau,\;729S^2D_g,\;S^2\delta C_y\bigr)=0.
\end{aligned}}
\end{equation}
It has $154$ operations and degree $784$. Both domains are unrestricted
integers. The file names \file{rational.json} and \file{integral.json}
refer to their auxiliary presentations, not to the domains of the solutions.

\section{Fixed constants and the inherited logical endpoint}
\label{sec:constants}

The six constants are the complete source34 tuple authenticated in the
pinned scalar-carrier packet \cite{scalar}. Three have short closed forms:
\[
 \beta=2^{1064921}-1,\qquad A_c=2^{1113},\qquad B_c=2^{666104}.
\]
The other three are supplied as complete, minimally encoded little-endian
nonnegative integers in the compressed payloads. Their bit lengths are
\begin{center}
\begin{tabular}{lrr}
\toprule
Constant & Bit length & Role\\\midrule
$\beta$&1064921&source base minus one\\
$K_0$&1065043&fixed source stencil\\
$M_C$&1064921&source mask\\
$M_F$&1064922&already shifted source mask\\
$A_c$&1114&determinant scale\\
$B_c$&666105&determinant scale\\\bottomrule
\end{tabular}
\end{center}
The authoritative binding file gives their byte counts and SHA-256 digests;
the verifier decompresses and authenticates every byte. These are fixed
numerals, not extra variables or consistency-dependent choices. No port
from a different, smaller coefficient tuple is substituted.

The inherited endpoint is the signed native/Pell source theorem for this
exact tuple. Its completed ordinary argument identifies source witnesses
with finite refutations, through the pinned native compiler and controller
calibration. The local proofs retained here are
\file{PRIME_TWELVE.md}, \file{INTERPOLATION12.md},
\file{signed-quadratic-proof.md}, \file{RESIDUE_AND_DETERMINANT.md}, and the
scalar-parent proof \cite{scalar,prime,interpolation}. The source theorem
is used as an explicit earlier theorem, just as a construction might use
a published Diophantine representation theorem. The current work does
not replace its ordinary proof by a new proof-assistant derivation.

The new contribution must therefore be read precisely: the equations are
fully instantiated, and the new unbounded reductions are proved below;
the logical interpretation composes those reductions with the cited
existing calibration. Finite circuit tests alone do not certify that
calibration and do not establish \eqref{eq:target} on their own.

\section{Why a weaker carrier is sufficient}

The old scalar carrier was designed to force $H\mid K$ directly. But its
later use is to show that a certain explicit integer $B$ vanishes by
combining divisibility with a size estimate. That last step admits a
weaker hypothesis.

\begin{lemma}[Power congruence with a height bound]\label{lem:heightbasic}
For nonzero integers $H$ and positive integers $d$, if
$H\mid B^d$ and $|B|^d<|H|$, then $B=0$.
\end{lemma}
\begin{proof}
The integer $B^d/H$ has absolute value less than one, so is zero.
\end{proof}
The lemma is elementary. Its useful role here is to change what the
upstream carrier must certify, while retaining the final exact conclusion.
This is a reduction in the equation, not an approximate replacement of a
logical condition.

Put
\begin{equation}\label{eq:pi}
 \pi(u)=(u^2+3)(u^3-19),\qquad
 \Phi_D(z)=D^5\pi(z/D)
 =z^5+3D^2z^3-19D^3z^2-57D^5.
\end{equation}
The retained intersectivity proof establishes a root of $\pi$ modulo every
positive integer, including powers of $2,3,19$. The polynomial has no
rational root: $u^2+3>0$ over the reals and $19$ is not a rational cube.

\begin{lemma}[Two congruences in one carrier]\label{lem:carrier}
Let $M,H$ be nonzero coprime integers, $D=MH$, and
\[
 C_*=MK_*+HL,\qquad z_*=DT+C_*.
\]
If $M^5H\mid\Phi_D(z_*)$, then $M\mid L$ and $H\mid K_*^5$.
\end{lemma}
\begin{proof}
For $p^e\Vert M$, suppose $b=v_p(z_*)<e$. In \eqref{eq:pi}, the leading
term has valuation $5b$. All other terms have valuations at least
$2e+3b$, $3e+2b$, and $5e$, strictly larger. Thus the sum has valuation
$5b<5e$, contradicting $M^5\mid\Phi_D(z_*)$. Hence $M\mid z_*$ and,
by coprimality, $M\mid L$. Reducing the carrier modulo $H$ gives
$H\mid(MK_*)^5$, and coprimality cancels $M$.
\end{proof}

The difference in strength is intentional: the $M$ part is recovered
exactly, while only a fifth-power congruence is required at $H$. This
allows the inexpensive source scale $Y=qv$ in \eqref{eq:base}; its
original cubic restriction is recovered from $M\mid v-q^2$.

\section{Soundness, with all exceptional branches in order}
\label{sec:soundness}

\subsection{Uniform arithmetic facts}
The actual $\beta$ is odd and is $1$ modulo $3$. Thus $q$ is even,
$q\equiv2\pmod3$, and $q\ne0$. Also $W\equiv4\pmod{12}$, so
\begin{equation}\label{eq:xlarge}
 |X|\ge4(\beta-1)>2^{36},\qquad X\equiv-1\pmod3.
\end{equation}
If $v=0$, then $a=0,H=\Delta=3$. Otherwise
$|a|=|qv(X+1)|>1000$. Therefore $\Delta\ge3$, $H$ is nonzero and odd,
and $Q$ is either zero or at least $9$. In particular $\delta$ never
vanishes. Since $a$ is divisible by $q$,
\[
 H\equiv3\pmod q,\qquad \gcd(H,M)=\gcd(H,2q^2)=1.
\]
Modulo $3$ one has
\begin{equation}\label{eq:mod3}
 a=H=\Delta=0,\quad D_f=D_\tau=1,\quad
 Q=0,\quad\delta=C_y=-1.
\end{equation}
For the first Pell radicand, $\ell=XY^2$ is $0$ or $-1$ modulo $3$.
Uniformly $D_f,D_\tau\ge1$, since $\ell(\ell+1)\ge0$ for every integer
$\ell$. No sign of $Y$ or $\ell$ has been assumed.

\subsection{Nonzero scale, square extraction and the strict guard}
If $S=0$, then $sr=0$ and hence $I=0$. Every scaled radicand vanishes;
both equations reduce to $A_*^{16}$. But
\[
 A_*=D^5\pi(T+C_*/D)\ne0
\]
because $D\ne0$ and the argument is rational. Thus every zero has
$s\ne0$ and $r\ne0$.

Divide a vanishing conjugate factor by $S$. For $P_{768}$ the four root
terms are
\[
 \sqrt{C_y/\delta},\quad3\sqrt{D_f},\quad9\sqrt{D_\tau},
 \quad27\sqrt{D_g}.
\]
For $P_{784}$ the first term is $\sqrt{\delta C_y}$. Their squared
$3$-adic orders are $0,2,4,6+v_3(D_g)$, omitting a zero guard term.
Within any nonempty signed subset of two or more terms there is a unique
least order, and it is even. Grouping by rational squareclass proves that
a rational signed sum forces every term rational: within a squareclass,
the unique least valuation cannot cancel. This is the retained
prime-order separation lemma, now applied before any positivity of $D_g$
is known. In particular $D_g$ is a nonnegative integer square.

Since $s\ne0$, $B_0\le0$ would give $D_g\le-3$. Therefore $B_0>0$.
The elementary joint-positivity lemma for
$xy-(t^2-x-y+1)^2$ gives
\begin{equation}\label{eq:strict}
 G>0,\qquad\eta\zeta>0.
\end{equation}
If $c=0$, then
$\eta\zeta=-Y(Y+1)k^2\le0$, a contradiction. Thus $c\ne0$, $I\ne0$,
and now $Q>1$. Adjacent-coefficient Pell descent makes the rational
auxiliary ordinate integral in
\begin{equation}\label{eq:aux}
 QV^2-(Q-1)y^2=1.
\end{equation}
In the integral presentation one first puts
$y=\sqrt{\delta C_y}/\delta$. The other three extracted roots are
integers $f,\tau,g$. This order handles the $c=0$ and $Q=0$ branches
without prematurely invoking a descent theorem that requires $Q>1$.

\subsection{Recover the original scale before the source bounds}
The integral signed norm relation gives $S\mid A_*$. Applying
Lemma~\ref{lem:carrier} with $L=v-q^2$ yields
\[
 2q^2\mid v-q^2,\qquad H\mid K_*^5.
\]
There is consequently an integer $u$ with
\begin{equation}\label{eq:recoverY}
 v=q^2(2u+1),\qquad Y=q^3(2u+1).
\end{equation}
The full original large-scale source front is restored here, before its
signed Pell estimates, rank extraction or native decoding is used. This
also excludes $v=0$ by an actual argument.

Let $K=X(X+2ac)-1$. The exact identity
\begin{equation}\label{eq:shortprojection}
 K_*=2K-HXc
\end{equation}
and oddness of $H$ give $H\mid K^5$. The short projector itself is
already present in the repository's positive-domain work; the new use
is its signed transfer into this weaker mixed-modulus carrier.

\subsection{Fold the main Pell root, then use the stronger height}
The retained signed main-folding argument uses the strong equation
$f^2-\Delta I^2=1$, the auxiliary equation \eqref{eq:aux}, the bound
$|c|\le|I|$ and $V\equiv-c\pmod f$. It gives
$|c|=\psi_{|a+2|}(p)$. The $c^2\mid I$ rank argument and the two
auxiliary residues then give $p=R$. These steps precede any main
projection. Their existing ratio bound supplies
\begin{equation}\label{eq:ratio}
 R=2r_0+1,\quad r_0\ge24,\quad x=|X|,\quad
 m=|Y|>\frac{(x-3)^{r_0}}2-1.
\end{equation}
With $\sigma=\sgn(c)$, the same projection identity gives
\[
 2^R K\equiv
 B:=(\sigma X-2^R)(2^R\sigma X+1)\pmod H.
\]
Thus $H\mid B^5$.

\begin{lemma}[Height at the actual coefficient tuple]\label{lem:height}
Under \eqref{eq:xlarge} and \eqref{eq:ratio}, $|B|^5<|H|$.
\end{lemma}
\begin{proof}
Each factor of $B$ has absolute value at most $2^{R+1}x$, whence
$|B|^5\le2^{20r_0+20}x^{10}$. Put $L=(x-3)^{r_0}$. Since $L\ge4$,
$m>L/4$, and
\[
 |H|\ge4(x-1)m-3>(x-1)L-3
 \ge\frac x2L\ge\left(\frac x2\right)^{r_0+1}.
\]
Finally
$x^{r_0-9}\ge2^{36r_0-324}>2^{21r_0+21}$ for $r_0\ge24$.
This is exactly the remaining comparison.
\end{proof}
Lemma~\ref{lem:heightbasic} gives $B=0$. Its second factor is odd and
nonzero, so $\sigma X=2^R$. The ordinary main projection and native
decoding now recover a finite refutation. This proves soundness. The
stronger estimate depends on the actual large $\beta$; it is not a
claim valid under only the earlier generic bound $x\ge64$.

\section{Completeness and the order of the witnesses}

Begin with a finite refutation and its canonical source data. Write the
old cubic-scale parameter as $u$, and set $v=q^2(2u+1)$. All quantities
depending on $Y$ are then unchanged. The canonical projection gives
$H\mid K$, hence $H\mid K_*$. Since $M\mid v-q^2$,
\[
 \kappa=C_*/D=K_*/H+(v-q^2)/M
\]
is an integer.

Choose the joint guard witness $t$ and its nonzero Pell witness $s$.
The retained full-residue interpolation theorem then chooses one strong
pair $I,f$ with $4sc^2\mid I$. Put $r=I/(sc^2)$. \emph{Fix this pair
before choosing the residue of $T$.} Intersectivity gives a root $u_0$
of $\pi$ modulo $I$. Request $T\equiv u_0-\kappa\pmod I$; interpolation
supplies such an auxiliary witness without changing $I$. Then
\begin{equation}\label{eq:completequotient}
 A_*=D^5\pi(T+\kappa),\qquad
 \frac{A_*}{S}=\frac{c^2H^4\pi(T+\kappa)}I\in\Z.
\end{equation}
For \eqref{eq:768} take
\[
 n=A_*/S+3f+9\tau+27g+y-V.
\]
For \eqref{eq:784} replace $y-V$ by $\delta(y-V)$. A conjugate norm
factor vanishes. Every coordinate is integral; $n$ has no positivity
restriction. The witness order is outer data, $t,s$, $I,f$, an intersective
residue, $T$, then $n$. In particular there is no circular choice of a
modulus constrained by a later residue.

\begin{theorem}\label{thm:12}
The two equations \eqref{eq:768} and \eqref{eq:784} have integer zeros
exactly when the pinned source has a finite refutation. Composing with
the existing source/calculus calibration proves \eqref{eq:target} for
both equations. Their exact degrees are $768$ and $784$.
\end{theorem}

\section{Centering, degree certificates, and the cheaper variants}

For $P_{768}$ let $U=Sn-A_*$, $P=SV$, $L=\delta P$, and
$N=S^2-P^2$. Pairing the two signs of the auxiliary radical gives
\begin{equation}\label{eq:rationalpair}
 \delta W_1^2+2LW_1+N,\qquad
 W_1=U-3Sf-9S\tau-27Sg.
\end{equation}
Taking the remaining three-root norm is exactly the cleared original
norm. For $P_{784}$ put $L=S\delta V$ and
$N=S^2\delta-L(SV)$; the paired expression is
\begin{equation}\label{eq:integralpair}
 W_1^2+2LW_1+N.
\end{equation}
These are universal ring identities. They remove the high terms before
degree propagation and correspond to the integral coordinate changes
$n\mapsto n+V$ or $n\mapsto n+\delta V$.

The relevant degrees are
\[
 \deg H=4,\ \deg M=2,\ \deg D=6,\ \deg I=4,
 \ \deg\delta=24,\ \deg V=21,\ \deg S=16,\ \deg A_*=35.
\]
In \eqref{eq:rationalpair}, $2LU$ has degree $96$, above the competing
$94$ and $74$. Its eight conjugates give degree $768$. In
\eqref{eq:integralpair}, $N$ has degree $98$, above the degree-$96$
mixed term, giving $784$. The exact nonzero coefficients
\begin{align}
 [j^{352}w^{96}v^{96}s^{48}r^{48}c^{80}T^{48}]P_{768}
 &=2^{744}3^{448}\beta^{352},\label{eq:leading768}\\
 [j^{320}w^{80}v^{80}s^{80}r^{80}c^{128}T^{16}]P_{784}
 &=2^{624}3^{400}\beta^{320}\label{eq:leading784}
\end{align}
prove attainment at the actual coefficients. The checker independently
derives the leading term at every live gate over all six formal ports;
the coefficients in \eqref{eq:leading768}--\eqref{eq:leading784} are
strictly positive after specialization.

For the cheaper variants, retain the original
$Y=(2v+1)q^3$, use $D=H$, $z_*=HT+K_*$ and $S=srH$, and use the same
$\Phi_H$. The guard-first $c\ne0$ argument is unchanged. The carrier
forces $H\mid K^5$, and Lemma~\ref{lem:height} resolves it. Without the
mixed-modulus circuitry, the integral and rational forms have respectively
degrees $1040,816$ in $140,144$ operations. Centering the integral form
reduces its degree to $800$ in $147$ operations. Its attained coefficient is
\[
 [j^{384}w^{96}v^{96}s^{48}r^{48}c^{80}T^{48}]P_{800}
 =2^{792}3^{480}\beta^{384}>0.
\]
Their complete proofs, graphs and checks are in \file{height12/};
the mixed-modulus forms are in \file{crt12/}.

\section{A second construction: combine signs of bounded intervals}

The direct branch starts from the pinned $75$-predicate raw C222 system
\cite{direct}. Its body has $1259$ coordinates and $1177$ quadratic
residuals. The first $594$ rows form a front encoding CRT, factorial,
binomial and flow data. Power requests in that front are certified by
shared Pell witnesses. The frozen source has an ordinary proof of its
finite-refutation equivalence, independent of K215 and the native tuple.

The direct logical endpoint is precise. The inherited
\file{KEndpoint.c222_iff_zf_refutation} theorem identifies
$C_{222}(1)$ with a finite proof of contradiction from ZF. The inherited
\file{ConstructibleChoice.zf_refutation_iff_zfc_refutation} theorem
identifies that with a finite ZFC refutation. The latter is the ordinary
constructible-universe equiconsistency theorem, without an added
consistency assumption. \file{Bridge.lean} defines \file{ConZFC} as the
negation of that exact ZFC proof predicate. These unconditional
calibration theorems, their original pins and the intervening ordinary
arithmetic proofs are included in \file{inherited/direct/} and the
immutable incoming archive. Their Lean import closure has not been
recompiled in this work. The later machine-specific endpoint in the
same upstream file is not needed by the direct construction.

\begin{lemma}[Shared signs]\label{lem:sign}
If $r+s=a+b$ and $a,b\ge0$, then
\[
 (r\ge0\ \wedge\ s\ge0)\quad\longleftrightarrow\quad rs\ge0.
\]
\end{lemma}
\begin{proof}
If $r<0$, then $s\ge-r>0$, so $rs<0$; interchange $r,s$ for the
other case. The forward implication is immediate. This includes zeros.
\end{proof}
For an integer $x$, Legendre's three-square theorem gives
\begin{equation}\label{eq:signrep}
 x\ge0\quad\longleftrightarrow\quad
 \exists a,b,c\in\Z\quad x=a^2+b^2+c(c+1).
\end{equation}
Indeed $4x+1$ is representable by three squares; modulo four exactly one
root is odd. Conversely $c(c+1)\ge0$ for every integer $c$. This is also
the explicit sign representation recorded in \cite{sun}.

Replacing two scalar sign conditions by one product condition therefore
uses one quadratic row and three auxiliary integers instead of two rows
and six. The actual seven replacements are shown below. Indices are
zero-based in the unchanged raw body. Each listed three-row range sums
exactly to $Q-r-s-1$.

\begin{center}\small
\begin{tabular}{llll}
\toprule
Product & Nonnegative anchor pair & Common $Q$ & Product / anchor rows\\
\midrule
$x_{562}x_{563}$ &$(x_{549},x_{550})$&$x_{547}$&537--539 / 526--528\\
$x_{585}x_{586}$ &$(x_{576},x_{577})$&$x_{574}$&553--555 / 547--549\\
$x_{593}x_{594}$ &$(x_{576},x_{577})$&$x_{574}$&559--561 / 547--549\\
$x_{601}x_{602}$ &$(x_{576},x_{577})$&$x_{574}$&565--567 / 547--549\\
$x_{619}x_{620}$ &$(x_{610},x_{611})$&$x_{608}$&578--580 / 572--574\\
$x_{637}x_{638}$ &$(x_{628},x_{629})$&$x_{626}$&591--593 / 585--587\\
$x_{123}x_{564}$ &$(x_{91},x_{537})$&$x_{530}$&540--542 / 518--520\\
\bottomrule
\end{tabular}\end{center}
The first six anchors retain both scalar guards. For the last,
rows $83$--$85$ give $x_{91}=x_{94}+x_{95}+1$, with $x_{94},x_{95},x_{537}$
still guarded. All fourteen deleted signs are thus recovered \emph{before}
any Pell, power or CRT semantics is invoked. The parent body assignment
is unchanged in both directions; only private sign witnesses are replaced.

This yields $68$ predicates. After complete integral monic projection it
gives $589$ variables, $371$ quadratic rows and $8157$ collected quartic
terms. The projected gain is $15$ variables, not the raw $21$-witness
saving: six product constraints prevent six formerly available coordinate
eliminations. This is why actual generated polynomials, rather than
subtraction of guard counts, are used for every size claim.

\section{Even indices remove the remaining small-ordinate guards}
\label{sec:even}

The next change is structural. Instead of certifying $m=n^k$, certify
\begin{equation}\label{eq:double}
 m^2=n^{2k}.
\end{equation}
The front already supplies $m\ge0$, $n\ge1$, and $k\ge1$, so \eqref{eq:double}
is equivalent to the original equality. Squaring an output in its
quadratic congruence row preserves degree two; doubling an index changes
only a scalar coefficient. The benefit is that every requested Pell
index is now even.

\subsection{Pell notation and a residue lemma}
For $a\ge2$ write
\[
 (a+\sqrt{a^2-1})^j=\chi_a(j)+\psi_a(j)\sqrt{a^2-1}.
\]
Both sequences are integral, $\chi_a$ is strictly increasing on
nonnegative indices, and $\psi_a(j)\ge j$. Standard reduction by the
fundamental norm-one unit $a+\sqrt{a^2-1}>1$ classifies every integral Pell solution,
up to independent signs, by these sequences. For a signed base $B$,
define $\chi_B,\psi_B$ by the same integer-polynomial recurrences.
For $b>0$ they satisfy $\chi_{-b}(j)=(-1)^j\chi_b(j)$ and
$\psi_{-b}(j)=(-1)^{j-1}\psi_b(j)$ for $j\ge1$.

\begin{lemma}[Abscissa residues]\label{lem:abscissa}
If $1\le i\le h$ and
$\chi_a(j)\equiv\pm\chi_a(i)\pmod{\chi_a(h)}$, then
$j\equiv\pm i\pmod{2h}$. In particular
$\chi_a(j)\equiv\pm a\pmod{\chi_a(h)}$ forces $j$ odd.
\end{lemma}
\begin{proof}
The doubling formulas give
$(a+\sqrt{a^2-1})^{2h}\equiv-1\pmod{\chi_a(h)}$ in the quadratic
ring. Reducing an index modulo $2h$ and reflecting about $h$ reduces
its abscissa to $\pm\chi_a(r)$ with $0\le r\le h$.
For $h\ge2$, the recurrence implies
$2\chi_a(h-1)<\chi_a(h)$. Two abscissae with indices below $h$
cannot sum to the modulus, so congruence up to sign forces equality,
hence $r=i$. If $i=h$, only the endpoint has zero residue. The case
$h=1$ is immediate. Taking $i=1$ proves the final assertion.
\end{proof}

\subsection{The new shared frame}
There are $21$ calls and $15$ distinct seed indices. Let $\mathcal O$
be the $21$ output coordinates and $\mathcal F$ the frozen set of $43$
distinct frame arguments. Define
\begin{equation}\label{eq:frame}
 W=4+2\sum_{v\in\mathcal F\setminus\mathcal O}v
       +2\sum_{m\in\mathcal O}m^2.
\end{equation}
Every summand is nonnegative before power interpretation. Consequently
$W\ge4$, and every call has $n<W$, $2k<W$, and $m^2<W$.
The unchanged frame equation is
\[
 A^2-W^3(W+2)Z^2=1.
\]
The retained first modulus guard forces $A\ge2$, hence $Z\ne0$.
Classifying this equation at base $W+1$, and reducing its ordinate
modulo $W$, shows the frame index is a positive multiple of $W$.
Thus
\begin{equation}\label{eq:framebound}
 A\ge(W+1)^W>W^W>W^2.
\end{equation}
For every base $n\ge1$, the unchanged modulus
$T_n=2An-n^2-1$ satisfies $T_n>A$.

\subsection{The parity argument, including arbitrary signs}
The common pair satisfies
\begin{equation}\label{eq:common}
 U^2-(A^2-1)V^2=1,\quad V\ge1,\quad
 V=\Bigl(\prod_{i=1}^{15}Y_i^2\Bigr)(v_0+1),\quad
 B=1+4V(q_0+1),\quad B\equiv A\pmod U.
\end{equation}
Here $v_0,q_0$ are the common-certificate coordinates, distinct from the
twelve-variable branch's notation. Every $Y_i$ is nonzero and $Y_i^2\mid V$
by the literal product identity. The odd $B$ cannot equal $\pm1$:
with $|U|=\chi_A(h)$, the residue $B\equiv A\pmod U$ excludes both
possibilities by size, including $h=1$.

For a seed with requested even index $K=2k$, retain
\begin{align}
 X^2-(A^2-1)Y^2&=1,& S_0^2-(B^2-1)R_0^2&=1,\\
 S_0&\equiv X\pmod U,&R_0&\equiv K\pmod{4Y}.
\end{align}
If $|Y|=1$, then $X=\pm A$. Writing the auxiliary solution with
arbitrary signs and transferring the base modulo $U$ gives
$\chi_A(j)\equiv\pm A\pmod{\chi_A(h)}$. Lemma~\ref{lem:abscissa}
forces $j$ odd. Since $B$ is odd, $\psi_B(j)\equiv j\pmod2$, so $R_0$
is odd. But its residue modulo $4Y$ forces it even. This contradiction
excludes $|Y|=1$.

Therefore the main index is at least two, and
\begin{equation}\label{eq:automaticbound}
 |Y|\ge\psi_A(2)=2A>K.
\end{equation}
The old index-slack nonnegativity guards are no longer needed to obtain
this strict bound. The argument allows negative $U,B,X,Y,S_0,R_0$.

\subsection{The index is exactly the requested one}
Suppose $|X|=\chi_A(i)$, $|Y|=\psi_A(i)$ and
$V=\psi_A(h)$. The elementary squared-divisibility identity says
\[
 \psi_A(i)^2\mid\psi_A(h)\quad\Longrightarrow\quad
 i\le h\ \text{ and }\ \psi_A(i)\mid h.
\]
For completeness, write $h=qi+r$ with $0\le r<i$. Reduction modulo
$\psi_A(i)$ and strict monotonicity give $r=0$; binomial expansion
modulo $\psi_A(i)^2$ then gives $\psi_A(i)\mid q$, hence the claim.

The abscissa residue gives $j\equiv\pm i\pmod{2h}$. Since
$B\equiv1\pmod{4|Y|}$, the ordinate residue gives
$j\equiv\pm K\pmod{4|Y|}$. As $|Y|\mid h$, these imply
$i\equiv\pm K\pmod{2|Y|}$. With
$1\le i\le|Y|$ and $1\le K<|Y|$, only $i=K$ is possible.

\subsection{Orient seed pairs before allowing base one}
There are designated supporting calls for all fifteen seeds, with
indices
\[
 0,1,5,7,9,10,12,13,14,15,16,17,18,19,20.
\]
Their bases are already at least two from the front, before any factorial
or power meaning is used. This qualification matters: the base
$x_{112}=1+x_{107}x_4^2$ of another call can initially be one.

For a supporting call put $e=2k$, $N=n^e$, $M_0=m^2$, and
$x=\chi_A(e),y=\psi_A(e)$. With
$H_0=x-(A-n)y$, $J_0=x+(A-n)y$, the recurrence gives
\[
 H_0\equiv N\pmod{T_n},\qquad
 H_0J_0=1+T_ny^2.
\]
The four possible signs of $(X,Y)$ in the modified power congruence
give respectively
\begin{center}
\begin{tabular}{ll}\toprule
Signs & Consequence\\\midrule
$(+,+)$ &$M_0\equiv N\pmod{T_n}$\\
$(-,-)$ &$T_n\mid M_0+N$\\
$(+,-)$ &$T_n\mid M_0N-1$\\
$(-,+)$ &$T_n\mid M_0N+1$\\\bottomrule
\end{tabular}\end{center}
The frame gives $M_0N<W^W<A<T_n$ and
$M_0+N<W+W^{W-1}<W^W$. The last three cases are impossible; in
particular $M_0N=1$ is impossible because $n\ge2,e\ge2$. Thus every
seed pair is positively oriented and $M_0=N$. Nonnegativity gives
$m=n^k$. After this orientation, the remaining calls, including those
with base one, use the positive pair and the simple bounded congruence;
no inverse-residue sign ambiguity remains.

\subsection{One derived pair needs two steps}
The old $t+1$ pair was computed by one Pell step from the $t$ seed.
After doubling, the seed index is $2t$ and the target is $2(t+1)$.
Thus two steps are essential. With stored coordinates for $AX$,
$(A^2-1)Y$ and $A^2$, the replacements are
\begin{align}
 X_{+2}&=2A\bigl(AX+(A^2-1)Y\bigr)-X,\\
 Y_{+2}&=2AX+(2A^2-1)Y.
\end{align}
They are quadratic in the stored coordinates, so no higher-degree
residual is introduced. Omitting this change would invalidate the compiler.

\section{Completeness and the literal quartic}

Given true front arithmetic data, compute the new $W$. Choose
\[
 A=\chi_{W+1}(W),\qquad Z=\psi_{W+1}(W)/W.
\]
The latter is integral by the recurrence modulo $W$. Choose positive main
pairs at every $2k$ and the derived pair by two steps. Let
$P=\prod_{i=1}^{15}Y_i^2$ and use the finite ring
\[
 (\Z/4P\Z)[u]/(u^2-(A^2-1)).
\]
The unit $A+u$ has finite multiplicative order. A positive multiple of
that order gives a common pair with $U\equiv1\pmod{4P}$ and
$V\equiv0\pmod{4P}$, so $V>0$, $P\mid V$ and $U$ is odd.
As $\gcd(U,4V)=1$, CRT gives $B\equiv A\pmod U$ and
$B\equiv1\pmod{4V}$. Choose $B>\max(A,1+4V)$ in that progression; monotonicity and the
recurrences give nonnegative comparison quotients. Choose the auxiliary
pairs at the same $2k$.
Polynomial dependence on the base gives both required congruences.

The power quotient is supplied by the exact recurrence identity
\begin{multline}
 \chi_A(e)-(A-n)\psi_A(e)-n^e\\
 =(2An-n^2-1)\sum_{j=0}^{e-1}n^{e-1-j}\psi_A(j).
\end{multline}
Thus all twenty-one modified power rows hold. This is a simultaneous
construction with the explicit order front, frame, main pairs, common
pair, CRT base, auxiliary pairs, quotients. No incompatible preselected
auxiliary witnesses are reused.

From a new zero, the preceding soundness proof recovers the original
twenty-one equalities at the same front tuple. Build a \emph{fresh old}
frame and certificate at indices $k$, using one step for its derived
$t+1$ pair, to obtain a zero of the pinned parent. Conversely, old
soundness supplies true front data, from which the new completeness
construction produces a zero. The maps concern zero existence, not
pointwise equality of old and new auxiliary coordinates.

The exact source edits consist of fifteen doubled seed residues,
twenty-one squared outputs, twenty-one squared frame contributions,
two derived-pair rows, and the doubled frame endpoint. Eleven scalar
index guards and one paired-index guard are deleted, leaving $56$
predicates after sign sharing. The complete source edit ledger and all
projection restorations are supplied in \file{direct566/}.

Let $q_1,\ldots,q_{372}$ be the literal quadratic residuals in its
\file{equation.json}, in the stated order of $566$ integer coordinates.
The second principal equation of this paper is
\begin{equation}\label{eq:566}
 \boxed{Q_{566}(x_0,\ldots,x_{565})=
              \sum_{i=1}^{372}q_i(x_0,\ldots,x_{565})^2=0.}
\end{equation}
The accompanying sparse coefficient file lists every one of the $6431$
nonzero monomials. All fixed coefficients are literal integers, and no
runtime call to a proof predicate or to a universal polynomial remains.
An integer zero of \eqref{eq:566} is exactly a common zero of those rows.

\begin{theorem}
The complete polynomial $Q_{566}$ has degree four, $566$ unrestricted
integer variables, $6431$ monomials, and coefficient height
$2445790974704$. Its integer-zero existence is equivalent to that of
the pinned direct parent and hence, using the existing direct
C222/calculus calibration, to finite ZFC refutability.
\end{theorem}

\section{Verification, limitations, and further research}

\subsection{What the checks establish}
The package includes independent checkers, ordinary proofs and full input
data. Verification is deliberately split into mathematical and finite
algebraic claims:
\begin{itemize}
\item The twelve-variable checkers authenticate all six payloads, prove
every named register by exact free-ring identities, verify the complete
$495$-term four-root norm and centered cuts, check acyclicity and liveness,
and derive attained leading coefficients through the actual DAGs.
\item The direct checkers bind the seven sign replacements to literal
source-row identities, bind every changed power/index/frame row, verify
the two-step correction and the actual fifteen-factor common product,
restore every projected coordinate by an integral polynomial, and expand
the entire final sum of squares independently.
\item Normal and optimized Python execute the same checks. Deliberately
changed projectors, carrier operations, source products, restorations
and coefficients are rejected. No mathematical acceptance depends on
Python assertions that disappear under optimization.
\item The degree and coefficient counts are exact. Finite Pell examples
and the negative control are supplementary diagnostics; the unbounded
proofs above establish the equivalences.
\end{itemize}

The supplied odd-index negative control shows why merely deleting the
index guards is invalid. It gives a genuine frame and a single generic
power certificate with $n=4,k=5,m=2$, $|Y|=1$, satisfying the squared-output
condition while $m\ne n^k$. The doubled-index residue rejects it. This
is a counterexample to a generic shortcut, not a root of the ZFC equation.

\subsection{Precise boundaries}
The words ``ZFC-provable'' describe an ordinary mathematical proof that
can be formalized inside ZFC, composed with the explicitly cited source
theorems. They do not claim that this package emits a new Lean/Rocq proof
of the entire native or direct compiler chain. The inherited calibration
and the new local algebraic certificates have different scopes and remain
separately identified. None of the work assumes ZFC consistency, and no
root search can establish the desired absence of roots.

The elementary ingredients are classical. The claimed progress is the
specific changes to the fully instantiated constructions at the pinned
snapshot, with exact new equations and verified comparisons. We do not
claim these are globally simplest, nor that the local lemmas have no
precedent anywhere in the literature.

\subsection{Further research questions}
\begin{enumerate}
\item Can more of the actual coefficient-size margin be exchanged for
weaker carrier divisibility, while preserving both directions of the
source reconstruction?
\item Can the $M^5$ condition be shortened by a proved restriction on the
prime factors of a canonical strong modulus? Generic modular roots do
not suffice: primes dividing the mandatory $c$ and $s$ must be covered.
\item The scale $Y=qv$ retains the uniform coprimality
$\gcd(H,2q^2)=1$. Dropping the last factor $q$ would reduce degree
further but breaks that uniform fact. Can a different carrier repair it
without a new variable or a circular coprimality assumption?
\item Can the even-index compiler be applied before the current raw
arithmetic expansion, so its source frame and shared certificates can be
redesigned together rather than locally rewritten?
\item Can integral changes of coordinates recover the six eliminations
lost when two scalar guards become a product guard? Can a larger row
basis expose additional independently oriented sign pairs?
\item Can the new transformations be formalized once as reusable
theorems, then applied to an authenticated complete native or direct
source theorem? The proof order around zero scale, base one, derived
indices and common periods should be explicit in such a formalization.
\end{enumerate}

\appendix
\section{A division-free norm evaluator}\label{app:norm}
This appendix makes explicit that the displayed norm is an ordinary
polynomial. Let the desired cleared norm be
\[
 d^8\NN(x;a,b,e,c/d).
\]
For the integral norm set $d=1$. Calculate in the quadratic algebra with
a formal symbol $u$ satisfying $u^2=e$. Put
\begin{align*}
 t_{20}&=x^2+e,&t_{21}&=-2x,&A&=a+b,\\
 q_0&=d(t_{20}-A)-c,&q_1&=dt_{21},\\
 T&=d^2ab,&R&=dc,\\
 v_0&=q_0^2+e q_1^2-4(T+RA),&v_1&=2q_0q_1,\\
 Z&=64TR,\\
 n_0&=v_0^2+e v_1^2-Zt_{20},&n_1&=2v_0v_1-Zt_{21}.
\end{align*}
The result is $n_0^2-e n_1^2$. Every operation is in the integer
polynomial ring. Expanding the product of all sixteen conjugates verifies
the identity with $495$ nonzero formal coefficients.

For the centered rational form, use $x=U+P$, $L=dP$,
$N=S^2-P^2$, and $c=dP^2-N$. Replace the expression for $q_0$ by
\[
 q_0=U(dU+2L)+N+d(e-a-b).
\]
For the centered integral form, use $x=U+L$, $c=L^2-N$, $d=1$, and
\[
 q_0=U(U+2L)+N+e-a-b.
\]
These local identities are checked separately before the full norm
identity is used. They explain the lower attained degrees without any
assumption that a residual has already vanished.

\section{Package map and reproduction}
The main deliverables are this PDF and its complete LaTeX source, together
with the following directories:
\begin{center}\small
\begin{tabular}{p{0.22\textwidth}p{0.69\textwidth}}
\toprule
Path & Contents\\\midrule
\file{crt12/}&Degrees $768/784$ literal DAGs, builder, independent checker,
complete coefficient payloads and full carrier proof.\\
\file{height12/}&Degrees $800/816/1040$ graphs and ordinary height/centering
proofs, with independent exact checks.\\
\file{direct566/}&Complete even-index quartic, quadratic residuals,
expanded monomials, source edits, projection and independent checks.\\
\file{direct589/}&The separately verified seven-sign-combination
intermediate and lower-height alternatives.\\
\file{peer/}&Independent mathematical and exact circuit reviews,
with the generic negative control in \file{direct566/}.\\
\file{inherited/}&Pinned ordinary proof dependencies, baseline review and
the original direct calibration archive.\\
\file{verify_all.py}&One entry point for all new checks and
deterministic circuit reproduction.\\
\file{SOURCE_LEDGER.json}&Pinned upstream identities, source locations,
and clear inherited-versus-new proof scope.\\
\bottomrule
\end{tabular}
\end{center}
From the extracted package root run
\begin{verbatim}
python3 verify_all.py
python3 -O verify_all.py
pdflatex -interaction=nonstopmode -halt-on-error report.tex
pdflatex -interaction=nonstopmode -halt-on-error report.tex
\end{verbatim}
The numerical verifiers use Python's standard library. The parent
coefficient recipes and full logical calibration are pinned to upstream;
the actual coefficients and all new final polynomials are included here.

\begin{thebibliography}{9}
\bibitem{repo} V. Reshetnikov, \emph{ProveIt}, snapshot
\file{84e16e43092614bae262ca59651cb9d77e31a255}, 24 October 2026.
\url{https://github.com/VladimirReshetnikov/ProveIt/commit/84e16e43092614bae262ca59651cb9d77e31a255}.
\bibitem{latest} ProveIt, \emph{Later arrivals and current size reductions},
\file{Computability/DiophantineConsistency/PROPOSAL_REVIEW_LATEST.md},
at the preceding snapshot. This is the baseline for the comparison table.
\bibitem{scalar} ProveIt, \emph{Scalar intersective carrier after signed
main-root recovery}, \file{Computability/DiophantineConsistency/candidates/scalar-carrier12/PROOF.md},
same snapshot; complete coefficient binding and graphs in that directory.
\bibitem{prime} ProveIt, \file{candidates/scalar-carrier12/source/PRIME_TWELVE.md},
within \file{Computability/DiophantineConsistency}, same snapshot.
\bibitem{interpolation} ProveIt,
\file{candidates/scalar-carrier12/source/INTERPOLATION12.md}, same snapshot.
\bibitem{direct} ProveIt,
\file{research/factorial-radius-output-sign-profile/raw/parent.json} and
\file{raw/domain.json}, within \file{Computability/DiophantineConsistency},
same snapshot; associated ordinary proofs and the active review identify
the direct C222, shared-CRT and Pell calibration.
\bibitem{compiler} ProveIt incoming package,
\file{ConZFC_407_888_Quartics_2026-10-24.zip}, same snapshot.
The unchanged terminal-elimination engine is used as a producer;
independent checkers validate its complete integral restoration certificates.
\bibitem{sun} Z.-W. Sun, \emph{Ten unknowns for Hilbert's tenth problem
over the integers}, arXiv:2609.01594v2, 2 September 2026.
\url{https://arxiv.org/abs/2609.01594v2}.
\end{thebibliography}
\end{document}
